\documentclass[12pt, fleqn]{article}

\usepackage[mathcal]{euscript}
\usepackage[all]{xy}

%\usepackage{graphicx}
\usepackage[top=30mm, bottom=40mm, left=25mm, right=25mm]{geometry}
\usepackage{amssymb}
\usepackage{epstopdf}
\usepackage{latexsym}
\usepackage[pdftex,bookmarks,breaklinks]{hyperref}
%\DeclareGraphicsRule{.tif}{png}{.png}{`convert #1 `dirname #1`/`basename #1 .tif`.png}

\usepackage{indentfirst}

\usepackage{marvosym}%Warning: this package screws up \Rightarrow and related commands

\begin{document}

\pagestyle{empty}
%\thispagestyle{empty}

\section*{What is ... Stuff?}

What is stuff made of?
A sensible answer to this question is a list of building blocks we can arrange in various ways to make the things we see around us. 
If this list is 
%not finite 
never-ending
then the question has no useful answer since we would not be able to compete such a list and therefore, there will always we stuff we cannot understand in the sense of the original question.
The Periodic Table of Elements on the wall of any chemistry class room is an example of such a finite list.

If elements on such a table cannot be further subdivided then we have found an answer to our question. 
This gives rise the notion of ``fundamental'' particles, that is, those that make up the final list on which no entry is further divisible. 

From here, we come to the notion of ``identical particles''. 
Any two of the same type on the list must be the same in every possible sense since, if they are not, there are really two such objects where we thought there was one.
In such a scenario, our original list was incomplete and we need to make a new list reflecting this distinction. 
Thus, if a finite final list exists, every entry is distinct and stuff made up of multiple copies of a particle in the list are identical particles of that type not distinguishable from one another in any fundamental sense. 

It is very difficult to deal directly with identical particles since they cannot be labeled (the label would count as a distinguishing feature and we could then extend our list again), but we can use a trick: We label them anyway by some randomly chosen assignment and then average over all ways we could have done this. 
For example, if we have two identical particles $XX$ of type $X$, say, we just call one of them $A$ and the other $B$: $XX \to AB$. Now consider switching the labelling $AB \to BA$. 
Doing this twice in a row does nothing to the labelling. 
There are two ways to ``average the labelling'' to reflect this fact: The symmetric average $AB+BA$ and anti-symmetric average $AB-BA$. 

%We have discovered that in order to treat the elements on the list as if we could label them ({\it i.e.}\ pretend identical particles are distinguishable) there is an ambiguity in doing so. 
%Surprisingly, when we compare this to the stuff we see around us, we find that only one assignment to each particle on the list we have so far is there! 
%We call the ones labeled symmetrically ``bosons'' and the ones labeled anti-symmetrically ``fermions''. 
%%There are two types we call ``boson'' and ``fermion''. The bosonic fundamental elements $X$ are those for with two copies $XX$ are labeled symmetrically whereas the fermionic ones are those for which the two copies are labeled anti-symmetrically.
%Apparently, from where we are sitting, the labelling symmetry is hidden or ``broken'' and much work at TAMU goes into figuring out how this happens in detail. 

We have discovered that in order to treat the elements on the list as if we could label them ({\it i.e.}\ pretend identical particles are distinguishable) there is an ambiguity in doing so. 
Surprisingly, when we compare this to the stuff we see around us, we find that each assignment can be thought of as a separate particle and these particles behave differently. 
We call the ones labeled symmetrically ``bosons'' and the ones labeled anti-symmetrically ``fermions'' and the symmetry exchanging them ``supersymmetry''.
The asymmetry between bosons and fermions is reflected in the mathematics we use to describe them: The bosons are most naturally described by the geometry of smooth objects and the fermions, by contrast, in terms of the algebra of discrete systems. 
These concepts are unified in ``superspace'' with motions in this space exchanging the two types of particles. 
Thus, the description of particles in this space automatically treats the bosons and fermions on the same footing. 

Apparently, from where we are sitting, this supersymmetry is hidden or ``broken'', and much work at TAMU goes into figuring out how this happens in detail. 
The question being addressed is how we are ``frozen in'' to our {\em specific} table of fundamental particles starting from the supersymmetric ultra-high-temperature epoch in the history of the Universe.

One of the paradigm-shifting lessons we are learning by studying this ``vacuum selection problem'' is that there seems to be even more symmetry than what we described as ``super''symmetry above. 
This symmetry exchanges not different different labellings of a particle but a particle together with a ``clump'' of particles of a different type! 
On the face of it, this is completely insane: The {\em fundamental} particles on our list are individually distinguishable by design but there are special {\em clumps} of some of these, that act like fundamental particles themselves. 
If such ``fundamental clumps'' are identical to fundamental particles already in our list, the description of this world is said to exhibit ``self-duality''. 

If this is not the case, the description of this world can be replaced with an equivalent description that looks like an entirely different world. In the language of the vacuum selection problem, these correspond to different vacua, but they are equivalent in the sense that the lumps of one are the fundamental particles in the list of the other and {\it vice versa}.

This suggests a radical change in our approach to the vacuum selection problem in which we look at all of the vacua related to each other by such dualities.
The proper mathematical definition of this so-called M-theory goes beyond what is available to date but certain properties of it are well-understood. 
In the TAMU theory group we are working to improve our understanding of this theory by leveraging the symmetries and dualities described in this article. 
In one approach, we are focusing developing a new list to replace the fundamental particles of a single vacuum with objects we call F-branes (working title) that describe the analogous thing for all M-theory vacua concurrently. 
Success in this approach would be the first formulation of M-theory that is ``background-independent'', that is, does not refer to any particular vacua for its definition. 

In a parallel line of investigation, we are translating what could be referred to as the most symmetric vacuum of M-theory into the language we use to describe the universe from our perspective. 
The strategy is to use the superspace appropriate to our particular list of fundamental particles as a lens through which to look at M-theory.
What we see when we put on our superspace glasses is strikingly simple: 
The theory is entirely determined by two mathematical structures that measure the volume of the space and the twisting of a surface called a ``membrane''. 
When a particular background is selected, the list of fundamental particles, their interactions, and how they move are all determined by these two quantities.

This encoding of M-theory ``dynamics'' in the geometry of superspace is elegant, but more importantly, it gives valuable insight into the process that freezes us into a specific vacuum as the Universe expands and cools. 
This is because, in this language, changes in the shapes of things as we see them in superspace are required to happen only in such a way that they preserve certain properties of the shape such as the volume and twisting mentioned above.
This conditions are very restrictive so this has direct consequences for the allowed paths to vacuum selection. 
In essence, the proposal is to chart these paths using our superspace as the surveying equipment. 

























\end{document}
